Bei der isothermen Expansion bleibt die innere Energie der Luft gleich. Nach dem ersten Hauptsatz wird die ganze am Lift verrichtete Arbeit als Wärme aus der Umgebung aufgenommen, $W = Q$. Wärme wird also vollständig in Arbeit umgewandelt.

Der zweite Hauptsatz (nach Planck) verbietet das aber nur für eine \emph{periodisch arbeitende} Maschine, die nach einem Zyklus wieder im Anfangszustand ist und als einzige Wirkung Wärme aus einem Reservoir in Arbeit umwandelt. Hier ist das nicht der Fall: Die Luft ist nach dem Vorgang nicht mehr im Anfangszustand, sie hat einen viel kleineren Druck und ein grösseres Volumen (ihre Entropie hat zugenommen). Diese Veränderung ist die zweite Wirkung des Prozesses.

Um den Vorgang zu wiederholen, müsste man die Luft wieder in die Flasche pressen. Dazu braucht ein Kompressor mindestens so viel Arbeit, wie beim Heben gewonnen wurde, und er gibt dabei die entsprechende Wärme wieder an die Umgebung ab. Die Pressluftflasche ist also ein Energiespeicher, der beim Füllen aufgeladen wurde, und kein Perpetuum mobile zweiter Art.
